\section{Revelation Theory and Algorithms}
\label{app:c}
Fix $f$, $\alpha\in(0,1)$, $\beta=1-\alpha$, and finite $\Gamma\geq1$. Use the
historical $\pi,p_1$, sensitivity bounds $[\underline p_\Gamma,\overline p_\Gamma]$,
and $s_\Gamma=\overline p_\Gamma-\underline p_\Gamma$. Write $\mu=P_X$ and
$\mathbb E=\mathbb E_\mu$. All kernels and policies are measurable; conditional
statements apply on components of positive population mass.

\subsection{Retained-information model and compatible laws}
\label{app:retained-class}\label{app:audit-information}
Additional revelation $A\in\{0,1\}$ satisfies $A\leq1-S$, with intensity
$\lambda(x)=P(A=1\mid X=x,S=0)\in[0,1]$. The audited population law has known
$p_A(x)=P(Y=1\mid X=x,A=1)$ satisfying
\begin{equation}
 \underline p_\Gamma\leq p_A\leq\overline p_\Gamma
 \quad\text{on }\{(1-\pi)\lambda>0\}.
 \label{eq:c-audit-bounds}
\end{equation}
Extend $p_A$ within these bounds on its zero-mass region; this has no effect on risks.

\begin{definition}[Retained-sensitivity compatibility]\label{def:c-retained-class}
The class $\mathcal P^{\mathrm{ret}}_{\Gamma,\lambda}$ consists of laws of
$(X,S,A,Y)$ with $X$-marginal $\mu$ and conditional components
\[
\begin{array}{c|c|c}
(S,A)&P(S,A\mid X=x)&P(Y=1\mid X=x,S,A)\\\hline
(1,0)&\pi(x)&p_1(x)\\
(0,1)&(1-\pi(x))\lambda(x)&p_A(x)\\
(0,0)&(1-\pi(x))(1-\lambda(x))&q(x).
\end{array}
\]
The residual kernel ranges freely over
$q\in[\underline p_\Gamma,\overline p_\Gamma]$; $(S,A)=(1,1)$ has zero mass.
\end{definition}

\begin{lemma}[Historical compatibility]\label{lem:c-compatible-laws}
Every admissible $q$ defines a retained law whose marginal belongs to
$\mathcal P_\Gamma(P_{\rm obs})$, with original hidden kernel
\begin{equation}
 p_0=\lambda p_A+(1-\lambda)q\quad\text{on }\{\pi<1\}.
 \label{eq:c-original-kernel}
\end{equation}
The retained class is nonempty.
\end{lemma}
\begin{proof}
Apply the kernel construction of Appendix~\ref{app:compatibility} to the table.
Its masses sum to one and its selected component is unchanged. Conditional total
probability gives \eqref{eq:c-original-kernel}; convexity of the sensitivity interval
proves compatibility. Taking $q=\underline p_\Gamma$ proves nonemptiness.
\end{proof}
At $\lambda=0$ the historical class is recovered; at $\lambda=1$ no residual
kernel remains. Membership requires the audited and residual bounds separately:
a bound on their mixture $p_0$ alone does not imply either component bound.

Known randomization $A\perp Y\mid(X,S=0)$ supplies extra information: $p_A=p_0$
on the audited support, and $q=p_0$ wherever residual mass remains. Thus the
population audited law identifies $p_0$ wherever originally hidden cases have
positive audit probability. The retained class instead keeps its residual kernel
free within the displayed bounds. Finite audits additionally require estimation
and uncertainty control for the audited contribution; a positive population audit
probability alone supplies no finite-sample precision guarantee.

\subsection{Post-revelation endpoints and tail identification value}
\label{app:tiv}
Use $a^y_{f,\eta},d_{f,\eta},b^\sigma_{f,\eta}$ from Section~\ref{sec:population},
write $s_\Gamma=\overline p_\Gamma-\underline p_\Gamma$, and set
$w_\lambda=(1-\pi)(1-\lambda)$. The known contribution and a residual
kernel's threshold objective are
\begin{align}
 K_{f,\eta}(\lambda)
 &=\mathbb E[\pi(a^0_{f,\eta}+p_1d_{f,\eta})
 +(1-\pi)\lambda(a^0_{f,\eta}+p_Ad_{f,\eta})],
 \label{eq:c-known-contribution}\\
 G_q(\eta)&=\eta+\beta^{-1}\{K_{f,\eta}(\lambda)
 +\mathbb E[w_\lambda(a^0_{f,\eta}+q d_{f,\eta})]\}.
 \label{eq:c-objective-q}
\end{align}
Conditioning on the table and optimizing the affine residual term give
\begin{equation}
 \begin{split}
 L_{f,\eta}(\lambda)&=\eta+\beta^{-1}
 \{K_{f,\eta}(\lambda)+\mathbb E[w_\lambda b^L_{f,\eta}]\},\\
 U_{f,\eta}(\lambda)&=\eta+\beta^{-1}
 \{K_{f,\eta}(\lambda)+\mathbb E[w_\lambda b^U_{f,\eta}]\}.
 \end{split}\label{eq:c-retained-objectives}
\end{equation}
\begin{proposition}[Sharp retained-sensitivity interval]\label{prop:c-sharp-endpoints}
The attainable population risks are exactly
\[
 \{\rho_{\alpha,Q}(f):Q\in\mathcal P^{\mathrm{ret}}_{\Gamma,\lambda}\}
 =[\min_{\eta\in[0,1]}L_{f,\eta}(\lambda),
   \min_{\eta\in[0,1]}U_{f,\eta}(\lambda)].
\]
Both threshold minima and endpoint laws are attained.
\end{proposition}
\begin{proof}
On $H=\{x:\pi(x)<1\}$, the common kernels $q_f^L,q_f^U$ of
Appendix~\ref{app:extremal} satisfy
$a^0_{f,\eta}+q_f^\sigma d_{f,\eta}=b^\sigma_{f,\eta}$ for every $\eta\in[0,1]$.
On $H^c$, the residual weight $w_\lambda=(1-\pi)(1-\lambda)$ is zero.
Thus $L_{f,\eta}(\lambda)\leq G_q(\eta)\leq U_{f,\eta}(\lambda)$ for every
admissible $q$, with equality at the corresponding common kernel.
Taking minima proves the endpoint claims;
continuity of each hinge objective on $[0,1]$ gives attainment. For
$q_t=(1-t)q_f^L+tq_f^U$,
\[
 \sup_\eta|G_{q_t}(\eta)-G_{q_s}(\eta)|
 \leq |t-s|\beta^{-1}\mathbb E[w_\lambda s_\Gamma].
\]
The infimum inequality \eqref{eq:a:infimum-stability} makes the optimized risk
continuous in $t$. This admissible path joins the endpoints and fills their interval.
\end{proof}

\begin{proposition}[Exact revelation value]\label{prop:c-tiv}
With
\begin{align}
 \tau_{f,\eta}&=(1-\pi)s_\Gamma|d_{f,\eta}|/\beta,
 \label{eq:c-tiv}\\
 W_{f,\eta}(\lambda)&:=U_{f,\eta}(\lambda)-L_{f,\eta}(\lambda)
 =\mathbb E[(1-\lambda)\tau_{f,\eta}]
 =W_{f,\eta}(0)-\mathbb E[\lambda\tau_{f,\eta}],
 \label{eq:c-linear-gap}
\end{align}
every bounded measurable $h$ and feasible $t$ satisfy
\begin{equation}
 W_{f,\eta}(\lambda+th)-W_{f,\eta}(\lambda)
 =-t\mathbb E[h\tau_{f,\eta}].\label{eq:c-feasible-derivative}
\end{equation}
\end{proposition}
\begin{proof}
Subtract \eqref{eq:c-retained-objectives}: $K$ cancels and
$b^U-b^L=s_\Gamma|d|$. Substitution gives the perturbation identity, including
one-sided feasible directions. Since $0\leq\tau\leq\beta^{-1}$, all terms are
integrable. The fixed-threshold reduction is independent of $p_A$, even when
that law varies with the policy.
\end{proof}

\subsection{Minimum-cost precision and the exact frontier}
\label{app:frontier}
Fix $f,\eta,\Gamma$, abbreviate $\tau=\tau_{f,\eta}$, and let $c\geq0$ be
measurable with $\mathbb E[c]<\infty$. Set $T=\mathbb E[\tau]=W_{f,\eta}(0)$,
$a=(T-\varepsilon)_+$ for $\varepsilon\geq0$, and
$F=\mathbb E[\tau\mathbf1_{\{c=0\}}]$. By \eqref{eq:c-linear-gap}, the
precision constraint is $\mathbb E[\tau\lambda]\geq a$.

\begin{proof}[Proof of Theorem~\ref{thm:cert-frontier}]
We establish attainment and the primal--dual identity
\begin{equation}
 \operatorname{Cert}_{f,\eta}(\varepsilon)
 =\min_{0\leq\lambda\leq1}\{\mathbb E[c\lambda]:\mathbb E[\tau\lambda]\geq a\}
 =\sup_{u\geq0}\{ua-\mathbb E[(u\tau-c)_+]\}.
 \label{eq:c-primal-dual}
\end{equation}
Pointwise, $(c-u\tau)\lambda\geq-(u\tau-c)_+$, so every feasible policy obeys
\begin{equation}
 \mathbb E[c\lambda]\geq u\mathbb E[\tau\lambda]-\mathbb E[(u\tau-c)_+]
 \geq ua-\mathbb E[(u\tau-c)_+].\label{eq:c-weak-duality}
\end{equation}
The following four cases attain this lower bound, or its limit.

\emph{No reduction ($a=0$).} Take $\lambda=0,u=0$, giving zero cost; this includes $T=0$.

\emph{Free reduction ($0<a\leq F$).} Take
$\lambda=(a/F)\mathbf1_{\{c=0,\tau>0\}}$ and $u=0$. The reduction is exactly
$a$ at zero cost, including when the free component has atoms.

\emph{Interior reduction ($F<a<T$).} On $D=\{c>0,\tau>0\}$, use the finite
measure $\nu(B)=\int_B\tau\,d\mu$ and $r=c/\tau$. Put $b=a-F\in(0,\nu(D))$
and choose
\begin{equation}
 u_* =\inf\{v\geq0:\nu(D\cap\{r\leq v\})\geq b\}.
 \label{eq:c-weighted-quantile}
\end{equation}
Because $r\in(0,\infty)$ $\nu$-almost everywhere, $0<u_*<\infty$.
Continuity of a finite measure along monotone sets gives
$H:=\nu(D\cap\{r<u_*\})\leq b\leq H+J$, where
$J:=\nu(D\cap\{r=u_*\})$. Set $\theta=(b-H)/J$ if $J>0$ and $\theta=0$
otherwise (then $H=b$). The policy
\begin{equation}
 \lambda_* =\mathbf1_{\{c=0,\tau>0\}}
 +\mathbf1_{D\cap\{r<u_*\}}+\theta\mathbf1_{D\cap\{r=u_*\}}
 \label{eq:c-optimal-policy}
\end{equation}
achieves reduction $F+H+\theta J=a$ and equality in the pointwise dual bound.
Its cost equals $u_*a-\mathbb E[(u_*\tau-c)_+]$. It ranks by $\tau/c$ with
threshold $1/u_*$; fractional ties handle atoms without an atomlessness assumption.

\emph{Full reduction ($a=T>F$).} Feasibility forces
$\mathbb E[\tau(1-\lambda)]=0$, hence $\lambda=1$ almost surely on $\{\tau>0\}$.
The minimum cost is $\mathbb E[c\mathbf1_{\{\tau>0\}}]$, attained by revealing
exactly that set, while
\[
 uT-\mathbb E[(u\tau-c)_+]=\mathbb E[\min\{u\tau,c\}]
 \uparrow\mathbb E[c\mathbf1_{\{\tau>0\}}]
\]
by monotone convergence. Thus the dual supremum agrees even without a finite
maximizer. Its shortfall is $\mathbb E[(c-u\tau)_+\mathbf1_{\{\tau>0\}}]$,
so a finite optimizer exists exactly when $c\leq u\tau$ almost surely on $\{\tau>0\}$ for
some finite $u$. Full and free reduction correspond to value-density thresholds
$0$ and $\infty$. Zero-value locations can always be left unselected; if $F=T$
the whole frontier is free. These cases exhaust $0\leq a\leq T$.
\end{proof}

\subsection{From threshold precision to optimized CVaR precision}
\label{app:optimized-width}
Denote the endpoints in Proposition~\ref{prop:c-sharp-endpoints} by
$\underline\rho^{\rm ret}_{\alpha,\Gamma,\lambda}(f)$ and
$\overline\rho^{\rm ret}_{\alpha,\Gamma,\lambda}(f)$, and their difference by
$I^{\rm ret}_{\Gamma,\lambda}(f)$.
\begin{proposition}[Controlling the optimized CVaR interval]\label{prop:c-optimized-width}
For the fixed predictor,
\begin{equation}
 0\leq I^{\rm ret}_{\Gamma,\lambda}(f)
 \leq\sup_{\eta\in[0,1]}W_{f,\eta}(\lambda)=W_{f,0}(\lambda).
 \label{eq:c-optimized-width}
\end{equation}
Theorem~\ref{thm:cert-frontier} at $\eta=0$ therefore minimizes cost for the
sufficient condition $W_{f,0}(\lambda)\leq\varepsilon$.
\end{proposition}
\begin{proof}
If $\eta_L$ minimizes $L$, then
$\min_\eta U-\min_\eta L\leq U_{f,\eta_L}-L_{f,\eta_L}=W_{f,\eta_L}$;
nonnegativity follows from ordered endpoints. For
$v=\max_y\ell(f(x),y)$ and $w=\min_y\ell(f(x),y)$,
\[
 |d_{f,\eta}(x)|=(v-\eta)_+-(w-\eta)_+=(v-\max\{w,\eta\})_+.
\]
This decreases in $\eta$ and equals $v-w$ at zero. Integration against
$(1-\lambda)(1-\pi)s_\Gamma/\beta\geq0$ proves the supremum identity.
\end{proof}

\paragraph{Direct optimized-width design.}
Fix an admissible set $\Lambda_{\rm adm}$ of measurable revelation intensities and
a specified family of population audit laws
$\{P^\lambda:\lambda\in\Lambda_{\rm adm}\}$ sharing one $(X,S,Y)$-marginal.
Each $P^\lambda$ has intensity $\lambda$ and satisfies
Definition~\ref{def:c-retained-class}. Set
$p_A^\lambda(x)=P^\lambda(Y=1\mid X=x,A=1)$, with an extension within the
sensitivity bounds on zero-mass regions.
Write $I^{\mathrm{ret}}_{\Gamma,\lambda}(f;p_A^\lambda)$ for the width obtained
using this kernel in \eqref{eq:c-known-contribution}, retaining the original
free residual bounds. For the cost density $c$ above and $\varepsilon\geq0$, define
\begin{equation}
 \operatorname{Cost}^{\mathrm{CVaR}}_f(\varepsilon)
 :=\inf_{\lambda\in\Lambda_{\rm adm}}
 \left\{\mathbb E[c\lambda]:
 I^{\mathrm{ret}}_{\Gamma,\lambda}(f;p_A^\lambda)\leq\varepsilon\right\},
 \label{eq:c-direct-optimized-cost}
\end{equation}
with $\inf\varnothing=+\infty$. The audited contribution can affect the separate
threshold minima even though it cancels from $W_{f,\eta}(\lambda)$.

\paragraph{Strictness of the sufficient bound.}
Take constant $X$, zero-one loss, $f=0$, $\pi=p_1=p_A=1/2$, $\lambda=1/2$,
$\Gamma=2$, and $\beta=1/4$. The original bounds are $[1/3,2/3]$.
Under the retained lower and upper endpoint laws, the probabilities $r_L$ and
$r_U$ of unit loss $L_f=1$ are
\[
 r_L=\tfrac14+\tfrac18+\tfrac1{12}=\tfrac{11}{24},\qquad
 r_U=\tfrac14+\tfrac18+\tfrac16=\tfrac{13}{24}.
\]
Both exceed $\beta$, so \eqref{eq:a:bernoulli-cvar} gives optimized endpoints
one and $I^{\rm ret}_{\Gamma,\lambda}(f)=0$, although
$W_{f,0}(\lambda)=\{(1/2)(1/2)/(1/4)\}(2/3-1/3)=1/3$.
Thus the sufficient criterion can require revelation when optimized risk is
already identified. Predictor reoptimization and estimated acquisition values
require their own statistical or sequential analysis.

\subsection{Finite-cell computation and operational cost units}
\label{app:finite-cell}
For cells $j=1,\ldots,K$ with masses $\mu_j>0$, write
$\tau_j=\tau_{f,\eta}(j)$, $c_j=c(j)$, and set
$v_j=\mu_j\tau_j$, $w_j=\mu_jc_j$, $T=\sum_jv_j$, $a=(T-\varepsilon)_+$,
and $F=\sum_{j:c_j=0}v_j$. The exact fractional algorithm is:
\begin{enumerate}
\item If $a=0$, reveal nothing. If $0<a\leq F$, reveal fraction $a/F$ of the
zero-cost, positive-value component; both cases cost zero.
\item Otherwise reveal all useful free cells. Order the positive-cost,
positive-value cells by decreasing $\tau_j/c_j$; set
$R_k=\sum_{j\leq k}v_j$, $C_k=\sum_{j\leq k}w_j$, and $R_0=C_0=0$.
\item For $b=a-F$, let $k=\min\{j:R_j\geq b\}$. Reveal the cells before $k$,
fraction $\theta=(b-R_{k-1})/v_k$ of cell $k$, and none after it. Then
\begin{equation}
 \operatorname{Cert}_{f,\eta}(\varepsilon)
 =C_{k-1}+\theta w_k
 =C_{k-1}+\frac{b-R_{k-1}}{\tau_k/c_k}.
 \label{eq:c-finite-frontier}
\end{equation}
\end{enumerate}
Tied ratios may be ordered arbitrarily; full reduction gives $\theta=1$ in the
last useful cell. If there are no positive-cost useful cells, the first step
covers every feasible target. This is Theorem~\ref{thm:cert-frontier}'s policy,
so the frontier is piecewise linear. For fixed cell policies and audited laws,
the piecewise-affine argument of Proposition~\ref{prop:b:threshold-search}
optimizes both retained objectives over
$\{0,1\}\cup\{\ell(f(j),y):1\leq j\leq K,\ y\in\{0,1\}\}$.

If a hidden label costs $\kappa(x)\geq0$, the target-population cost density and
operational value density are
\begin{align}
 c(x)&=(1-\pi(x))\kappa(x),\qquad
 \mathbb E[c\lambda]=\mathbb E[(1-\pi)\lambda\kappa],\label{eq:c-label-cost}\\
 \frac{\tau_{f,\eta}(x)}{c(x)}
 &=\frac{s_\Gamma(x)|d_{f,\eta}(x)|}{\beta\kappa(x)}
 \quad\text{when }(1-\pi)\kappa>0.\label{eq:c-operational-density}
\end{align}
The hidden-mass factor cancels; unit label cost ranks already-hidden cases by
$s_\Gamma|d|/\beta$. Cell masses likewise cancel from $v_j/w_j$.
The algorithm allocates divisible population mass. Whole-individual batches add
integer constraints and sampling variation. Sequential retraining changes losses
and acquisition values; confidence guarantees across adaptive rounds additionally
require control of newly revealed contributions, acquisition history, and stopping.
